\subsection{Resposta ordinal e comparações de interesse}
A resposta original, $s$, é o escore inteiro de popularidade. A pergunta preditiva é qual de três níveis a faixa ocupa. Agrupar o escore permite distinguir uma troca entre classes vizinhas de um erro entre os extremos. Também admite que escores diferentes, dentro do mesmo intervalo, correspondam à mesma decisão. A regra de classificação foi fixada no protocolo antes da avaliação das 9.526 faixas
\begin{equation}
Y(s)=
\begin{cases}
0, & 0\leq s<20 \quad \text{(baixa)},\\
1, & 20\leq s<40 \quad \text{(média)},\\
2, & 40\leq s\leq 100 \quad \text{(alta)}.
\end{cases}
\end{equation}
Os cortes reagrupam as cinco bandas do planejamento. Manter o limite inferior da classe alta em 40 dá suporte à avaliação dessa categoria em uma distribuição concentrada nos escores baixos. O estrato de 80--100, por exemplo, tem apenas cinco faixas na coorte. Os limites são uma definição operacional, sem interpretação oficial de sucesso musical. A Figura~\ref{fig:niveis} apresenta o escore inteiro e os cortes em cada conjunto. A composição da classe alta pode variar dentro de 40--100, mesmo quando o total permanece fixo.

\figura[.34\textheight]{popularidade_sorteio_coorte}{Distribuição do escore nas duas seleções de 9.526 faixas e cortes da classificação.}{fig:niveis}{Barras em porcentagem de cada conjunto, com peso igual por faixa.}

\begin{table}[H]\centering\small
\caption{Classes observadas nas 9.526 faixas da avaliação.}\label{tab:niveis}
\begin{tabular}{llrr}\toprule
Nível & Escore & Faixas & Percentual (\%) \\\midrule
Baixa & 0--19 & 5.548 & 58,24 \\
Média & 20--39 & 2.919 & 30,64 \\
Alta & 40--100 & 1.059 & 11,12 \\\midrule
Total &  & 9.526 & 100,00 \\\bottomrule
\end{tabular}
\end{table}

O Baseline tem 15 atributos de áudio, música e metadados. São 11 numéricos, ano, dançabilidade, energia, intensidade sonora, presença de fala, caráter acústico, instrumentalidade, indício de execução ao vivo, valência, andamento e duração. Os quatro categóricos são gênero, tonalidade, modo e assinatura de tempo. São as variáveis definidas e descritas na Seção~\ref{sec:fontes}. Títulos, créditos e IDs servem para integração e agrupamento. Idioma, peso amostral e informações de vinculação ficam fora do vetor de preditores.

\begin{table}[H]\centering\small
\caption{Quatro configurações e informação textual acrescentada ao Baseline.}\label{tab:blocos}
\begin{tabular}{@{}lP{.61\linewidth}r@{}}\toprule
Modelo & Atributos & Campos iniciais \\\midrule
Baseline & Áudio, música e metadados & 15 \\
Tradicional & Baseline e 21 descritores lexicais e estruturais & 36 \\
Profundo & Baseline e 384 coordenadas de embeddings & 399 \\
Combinado & Baseline, 21 descritores e 384 coordenadas & 420 \\\bottomrule
\end{tabular}\nota{Contagens antes da expansão das variáveis categóricas.}
\end{table}

As quatro configurações correspondem a perguntas de desempenho sobre os mesmos exemplos. O Baseline informa o que áudio e metadados já permitem classificar. Tradicional versus Baseline mede o incremento dos descritores. Profundo versus Baseline mede o incremento dos embeddings. Combinado versus Baseline mede o ganho das duas representações juntas, e Combinado versus Profundo examina se os descritores acrescentam informação aos embeddings. Para cada par, a expectativa de melhoria corresponde a uma diferença positiva na acurácia balanceada. Os contrastes são avaliados de forma pareada e interpretados com as outras métricas, sem exigir uma ordenação prévia dos resultados.

\subsection{Normalização e descritores lexicais e estruturais}
A letra elegível é convertida ao mesmo formato textual para os dois blocos. Aplicam-se normalização Unicode NFC, padronização das quebras de linha, remoção de espaços finais e remoção de linhas vazias externas. Palavras, diacríticos, sinais internos, ordem de linhas e repetições são preservados. A versão normalizada serve para produzir atributos, enquanto texto e metadados de origem permanecem documentados. A segmentação não depende de um dicionário de um idioma específico.

A motivação dos descritores é comparar propriedades mensuráveis de extensão, diversidade lexical, repetição e composição gráfica. Uma letra longa pode repetir muitas vezes o mesmo refrão. Por isso, são consideradas tanto contagens de tamanho quanto medidas relativas de diversidade e repetição. A Tabela~\ref{tab:descritores} explicita como cada variável foi obtida.

As medidas são calculadas no texto normalizado comum. Denotam-se por $C$ os caracteres não brancos, por $L$ as letras Unicode, por $T$ os segmentos lexicais, por $U$ os tipos distintos e por $B$ as linhas não vazias. Segmentos são sequências de letras que admitem apóstrofo ou hífen interno, após uniformização de caixa por \textit{casefold}. Razões com denominador nulo recebem zero. O indicador 21 registra a ausência de segmentos.

\begingroup\footnotesize\singlespacing
\begin{longtable}{rP{.31\linewidth}P{.56\linewidth}}
\caption{Definições operacionais dos 21 descritores de letras.}\label{tab:descritores}\\
\toprule & Descritor & Definição \\\midrule\endfirsthead
\toprule & Descritor & Definição (continuação) \\\midrule\endhead
\bottomrule\endfoot
\bottomrule\endlastfoot
1 & Caracteres totais & Pontos de código, incluindo espaços e quebras internas.\\
2 & Caracteres não brancos & Contagem $C$.\\
3 & Linhas não vazias & Contagem $B$.\\
4 & Estrofes operacionais & Blocos não vazios separados por uma ou mais linhas vazias.\\
5 & Segmentos lexicais & Contagem $T$.\\
6 & Tipos lexicais & Contagem $U$, após casefold.\\
7 & Comprimento médio de segmento & Média de caracteres dos segmentos após casefold.\\
8 & Dispersão de comprimento de segmento & Desvio com divisor $T$.\\
9 & Razão tipo/segmento & $U/T$.\\
10 & Razão de tipos únicos em frequência & Número de tipos com uma ocorrência dividido por $U$.\\
11 & Entropia de segmentos & $-\sum_{t=1}^{U}p_t\log_2p_t$, com frequências relativas $p_t$.\\
12 & Comprimento médio de linha & Média de caracteres não brancos nas linhas não vazias.\\
13 & Dispersão de comprimento de linha & Desvio com divisor $B$.\\
14 & Fração de linhas repetidas & $(B-\text{número de linhas distintas})/B$, após casefold e uniformização de espaços.\\
15 & Fração de letras & $L/C$, com categoria Unicode L.\\
16 & Fração de dígitos & Dígitos decimais Unicode divididos por $C$.\\
17 & Fração de pontuação & Caracteres de categoria Unicode P divididos por $C$.\\
18 & Fração de maiúsculas & Letras maiúsculas divididas por $L$.\\
19 & Linhas com interrogação & Fração das $B$ linhas com interrogação simples ou de largura completa.\\
20 & Linhas com exclamação & Fração das $B$ linhas com exclamação simples ou de largura completa.\\
21 & Ausência de segmentos & Indicador de $T=0$.\\
\end{longtable}
\endgroup

A entropia segue a definição discreta em bits de \citeonline{cover2006}, aplicada aqui aos segmentos. Ela e a razão tipo/segmento dependem da extensão e da segmentação. O indicador 21 permite registrar a ausência de segmentos na regra adotada.


\subsection{Embeddings multilíngues}
O \textit{multilingual-e5-small} é um encoder treinado por aprendizagem contrastiva em múltiplos idiomas. Ele fornece um vetor denso de 384 coordenadas a partir do texto \cite{wang2024,e5card}. A escolha permite representar informação contextual além de contagens lexicais e usar a mesma transformação para os idiomas presentes. O encoder é congelado na revisão \texttt{614241f6}. O modelo de popularidade ajusta coeficientes sobre os vetores, sem atualizar os parâmetros da rede.

O texto completo é tokenizado e repartido em janelas de até 512 posições, incluindo tokens especiais e o prefixo \texttt{query: }. Janelas consecutivas usam sobreposição de 32 tokens de conteúdo. Esse procedimento incorpora trechos além da primeira janela, enquanto mantém algum contexto na transição. Para a janela $j$, denote-se por $\bm h_{jt}$ o estado do token $t$, por $a_{jt}$ sua máscara de atenção e por $c_j$ a quantidade de tokens de conteúdo novos. Calculam-se a média mascarada, a normalização da janela e a normalização da soma ponderada
\begin{equation}
\bm v_j=\frac{\sum_t a_{jt}\bm h_{jt}}{\sum_t a_{jt}},\qquad
\bm u_j=\frac{\bm v_j}{\|\bm v_j\|_2},\qquad
\bm z=\frac{\sum_j c_j\bm u_j}{\left\|\sum_j c_j\bm u_j\right\|_2}.
\end{equation}
A máscara exclui o preenchimento. Os pesos reduzem a influência de contar novamente posições de sobreposição, embora a rede possa usar esse contexto em mais de uma janela. A média mascarada e a normalização seguem a documentação do encoder. A agregação por janelas é uma adaptação do estudo.

A tokenização separada de prefixo e conteúdo introduz um token de espaço adicional em comparação com a concatenação textual do exemplo oficial. Essa diferença foi identificada na auditoria da implementação. O experimento avalia a variante efetivamente executada, com a mesma configuração para toda a coorte. Houve 9.500 entradas textuais distintas para o encoder. Foram reutilizados 9.018 vetores cuja entrada e configuração coincidiam com o registro anterior, e 482 foram calculados para completar a avaliação.

\subsection{Regressão logística ordinal regularizada}
A resposta tem ordem e duas fronteiras. O modelo de chances proporcionais representa a ordenação por dois limiares, usando o mesmo vetor de coeficientes para ambas as fronteiras \cite{mccullagh1980,agresti2010}. A restrição é parcimoniosa diante de um bloco de 384 coordenadas e de atributos correlacionados. A penalização quadrática ajuda a estabilizar o ajuste nessa dimensão \cite{hoerl1970,hastie2009}.

Para o vetor de preditores $x_i$, define-se
\begin{equation}
F_k(x_i)=\Pr(Y_i\leq k\mid x_i)
=\frac{1}{1+\exp[-(\theta_k-x_i^\top\beta)]},\qquad
k=0,1,\quad \theta_0<\theta_1.
\label{eq:logit}
\end{equation}
Então $p_{i0}=F_0(x_i)$, $p_{i1}=F_1(x_i)-F_0(x_i)$ e $p_{i2}=1-F_1(x_i)$. A Equação~\ref{eq:logit} supõe que uma mudança em $x_i^\top\beta$ desloca as duas chances acumuladas da mesma maneira na escala logit. Para um conjunto de treinamento $\mathcal T$ com $n_t$ faixas, minimiza-se
\begin{equation}
-\frac{1}{n_t}\sum_{i\in\mathcal T}\log p_{i,y_i}
+\frac{\alpha}{2}\sum_j \beta_j^2.
\end{equation}
Os limiares não são penalizados. Os coeficientes recebem a penalização L2 com intensidade $\alpha$. As faixas recebem pesos iguais no ajuste. Os limiares são parametrizados com distância positiva, e a otimização usa L-BFGS-B com gradiente analítico.

\subsection{Validação agrupada e seleção do hiperparâmetro}
Faixas de um mesmo artista, participações explícitas, composições relacionadas pela chave operacional e letras iguais ligam observações entre si. Construiu-se um grafo com essas ligações. Cada componente conexo permanece na mesma parte para impedir que relações reconhecidas atravessem a separação entre treino e teste. Foram obtidos 6.271 componentes. O agrupamento usa metadados e identidade textual, sem usar o escore para distribuir partes. O algoritmo procurou equilibrar o número de faixas, e as três classes tiveram suporte nas partes de treino e avaliação.

A validação cruzada tem cinco partes externas e três internas. Em cada iteração externa, uma parte é reservada para avaliação. Dentro das outras quatro partes, as três divisões internas comparam $\alpha\in\{0{,}001,0{,}01,0{,}1,1\}$ pela acurácia balanceada reunida das previsões internas. Empates favorecem a penalização maior. Imputação pela mediana, padronização e níveis da codificação indicadora são ajustados no respectivo treino, com tratamento de categorias inéditas. A combinação selecionada é ajustada novamente ao treino externo inteiro e avaliada na parte reservada \cite{cawley2010}. Cada faixa tem uma única previsão externa por modelo. Foram usados 240 ajustes internos e 20 avaliações externas para os quatro blocos.

A classe prevista é a de maior probabilidade, com empate resolvido pela classe de menor índice. Uma referência adicional atribui as frequências das classes no treino externo como probabilidades e usa a mesma regra de decisão. Isso dá uma medida de quanto se pode acertar com a prevalência das classes sozinha.

\subsection{Métricas e estabilidade dos contrastes}
Para $C_{kj}$ igual ao número de faixas de classe observada $k$ e prevista $j$, a sensibilidade é $r_k=C_{kk}/\sum_j C_{kj}$ e a precisão é $q_k=C_{kk}/\sum_j C_{jk}$, quando o respectivo denominador é positivo. A sensibilidade mede a fração de uma classe observada que foi recuperada. A precisão indica a fração correta entre as previsões daquela classe. Reunindo as previsões externas por faixa, calculam-se
\begin{equation}
\mathrm{AB}=\tfrac13\sum_{k=0}^{2}r_k,\qquad
\mathrm{Ac}=\tfrac1n\sum_{k=0}^2 C_{kk},\qquad
\mathrm{F1}_{\mathrm{macro}}=\tfrac13\sum_{k=0}^2\frac{2C_{kk}}{\sum_j C_{kj}+\sum_j C_{jk}}.
\end{equation}
A acurácia balanceada (AB) é a métrica principal porque atribui a mesma importância à recuperação de cada classe, mesmo quando seus tamanhos são diferentes \cite{brodersen2010}. A acurácia comum representa a fração de acertos e reflete as prevalências. O macro-F1 incorpora precisão e sensibilidade. A taxa de trocas entre os níveis baixo e alto é $(C_{02}+C_{20})/n$. Para as probabilidades externas $\widehat p_{ik}$, as medidas são \cite{gneiting2007}
\begin{equation}
\operatorname{LogLoss}=-\tfrac1n\sum_i\log\widehat p_{i,y_i},
\qquad
\operatorname{Brier}=\tfrac1n\sum_i\sum_{k=0}^2(\widehat p_{ik}-\mathbf{1}\{y_i=k\})^2.
\end{equation}
Valores menores indicam probabilidades mais adequadas segundo cada critério. A calibração descritiva compara a média prevista e a frequência observada em dez intervalos de largura 0,1 por classe. Na avaliação ponderada, os pesos $w_i^L$ são aplicados aos mesmos acertos e erros. A comparação principal conserva pesos iguais por faixa.

A estabilidade de $\Delta\mathrm{AB}$ e da diferença de sensibilidade da classe alta foi examinada com bootstrap pareado por componentes nas previsões externas congeladas \cite{field2007}. Em cada uma de 8.000 réplicas (semente 42), sorteiam-se 6.271 componentes com reposição. Um componente sorteado leva todas as suas faixas e as previsões dos quatro modelos. Para multiplicidades $b_g$, o contraste de $A$ versus $B$ é
\begin{equation}
\Delta\mathrm{AB}^{*}=\frac13\sum_{k=0}^2
\frac{\sum_g b_g (D^A_{gk}-D^B_{gk})}{\sum_g b_g N_{gk}},
\end{equation}
em que $D^A_{gk}$ conta os acertos de $A$ e $N_{gk}$ as faixas da classe $k$ no componente $g$. Para a sensibilidade alta, retém-se apenas o termo $k=2$ sem o fator $1/3$. Usam-se os percentis 2,5 e 97,5, separadamente para cada contraste. Essa análise quantifica a variação ao mudar a composição dos componentes avaliados, com previsões fixas. A variabilidade integral incluiria a seleção da coorte, as divisões, a escolha do hiperparâmetro e novos ajustes. As cinco avaliações externas compartilham observações de treino e, portanto, são dependentes \cite{bengio2004,bates2024}.

